\subsubsection{证据分层与继续研究的验收条件}
本文把四类内容分开：Dette--Titoff 原文给模型判别优化框架；软历史代理、方向与弦界、共享平方压缩、AC 逆映射及弱边先验界为自行推导；程序实现和求解状态由源码与冻结配置核查；排错、参数误差、误拒绝及耗时来自独立历史评价。
局部求解取得的似然没有完整 AC 全局最优证书。代理界仅在同一冻结矩阵和解释域内有效；一点 AC 回代不足以证明全域非线性误差界。

失败也须按原因区分。弱支路解析界给一类声明先验及设备预算下的信息上限；全局阻抗反例显示局部满秩仍不足以保证参数唯一。
预测失配与代理高分而完整 AC 竞争解释仍接近，属于模型及近似范围问题；原上界超时和连续积分退化属于计算问题。
一条记录可能同时含多个原因，目前没有证据把总体净差逐例唯一归因。

下一轮先以同判据检验共享平方 MILP 的求解精度和预算，再重新冻结策略及整历史校准，不能用已评价历史调策略后继续称独立评价。
滚动与固定首轮计划须在共同历史、同电量与共同外生未来下完成独立配对校准；耗时 gate 未过时保留停止记录。
proper EIG 先逐成分通过历史证据、独立重复和 latent 倍增检查，再做少量动作的预测积分与先验敏感性。显示质量小不能成为删除难积分成分的理由。
预测权重尚未校准，统一 AC 代理误差尚未建立，弱支路仍未决；当前保持三末端，不强行扩网。
原八小时协议与预算后用户续接分别保存，完整未完成事项和直接运行入口见 \path{adaptive_resume_plan.md} 与 \path{adaptive_milp_checkpoint.json}。
